% CLAUDE REPORT::WCDERIV-0001::2026-11-11T05:50-08:00
\documentclass[11pt]{article}
\usepackage{amsmath,amssymb,amsthm}
\usepackage[margin=1in]{geometry}
\newtheorem{theorem}{Theorem}
\newtheorem{corollary}[theorem]{Corollary}
\theoremstyle{definition}
\newtheorem{definition}[theorem]{Definition}
\newcommand{\WR}{\overline W_R}
\title{Worst-Case Minimum Derivation Cost and the Bridge\\
$\mathrm{SAT}\in\mathrm{P}\Rightarrow \overline W_R(N)\le\mathrm{poly}(N)$\\
\large WCDERIV-0001 (scheduled Claude run)}
\date{2026-11-11 05:50 America/Los\_Angeles}
\begin{document}
\maketitle

\paragraph{Status.} P versus NP remains \textbf{OPEN}. This run proves
elementary but load-bearing facts about the shared target. The main one is that,
for every proof system where the lower-bound half is already known, the bridge
is \emph{equivalent} to $\mathrm{P}\ne\mathrm{NP}$. Labels: PROVED/EXACT,
KNOWN THEOREM, COMPUTATIONAL, OPEN.

\section{Objects}
\begin{definition}
Let $R$ be a Cook--Reckhow refutation system: a polynomial-time verifier $V_R(F,\pi)$
such that $F\in\mathrm{UNSAT}$ iff $\exists\pi\,V_R(F,\pi)=1$. Take
$\operatorname{cost}_R(\pi)$ polynomially related to $|\pi|$, so that it charges
every step and every stored intermediate state. Define
\[
\WR(N)=\max_{F\in\mathrm{UNSAT},\ |\langle F\rangle|\le N}\ \min_{\pi:F\vdash_R\bot}\operatorname{cost}_R(\pi),
\]
together with $\mathrm{LB}_R$: ``$\WR$ is superpolynomial'' and
$\mathrm{Bridge}_R$: ``$\mathrm{SAT}\in\mathrm{P}\Rightarrow\WR(N)\le\mathrm{poly}(N)$''.
\end{definition}
The multisign DAG-hypergraph reduction system $R_{\mathrm{hyp}}$ fits this
definition whenever each reduction rule is sound and can be checked locally in
polynomial time. Its ``minimal work to reach the minimal reduced state'' is then exactly
$\min_\pi\operatorname{cost}_{R_{\mathrm{hyp}}}(\pi)$. Merging signed hyperedges on a
complementary literal is resolution, so the facts below apply to $R_{\mathrm{hyp}}$ directly.

\section{Results}
\begin{theorem}[PROVED/EXACT: the bridge is conditional p-optimality]\label{t:opt}
For every Cook--Reckhow system $R$:
$\mathrm{Bridge}_R\iff\big(\mathrm{SAT}\in\mathrm{P}\Rightarrow R\text{ is polynomially optimal}\big)$.
\end{theorem}
\begin{proof}
Suppose a polynomial-time algorithm $A$ decides SAT. Then $S_A$, which accepts the empty
proof of $F$ iff $A(F)=0$, is a polynomially bounded system. ($\Leftarrow$) If $R$
simulates $S_A$ with polynomial blowup, then $\WR$ is polynomially bounded.
($\Rightarrow$) A polynomially bounded $R$ simulates every system, trivially.
\end{proof}

\begin{theorem}[PROVED/EXACT: bridge collapse]\label{t:collapse}
If $\mathrm{LB}_R$ holds unconditionally, then $\mathrm{Bridge}_R\iff\mathrm{P}\ne\mathrm{NP}$.
\end{theorem}
\begin{proof}
$\mathrm{Bridge}_R\wedge\mathrm{LB}_R\Rightarrow\mathrm{SAT}\notin\mathrm{P}$.
Conversely, if $\mathrm{P}\neq\mathrm{NP}$ then $\mathrm{Bridge}_R$ holds vacuously.
\end{proof}
$\mathrm{LB}_R$ is a KNOWN THEOREM for resolution (Haken 1985; Ben-Sasson--Wigderson width),
tree-like resolution, cutting planes (Pudl\'ak 1997), bounded-depth Frege
(Ajtai; Beame--Impagliazzo--Pitassi; Kraj\'{\i}\v{c}ek--Pudl\'ak--Woods),
and polynomial calculus (Razborov). For each of these systems, and for any
$R_{\mathrm{hyp}}$ that is at most as strong as resolution, the ``weakest sufficient bridge''
is therefore \emph{the whole problem}.

\begin{corollary}[PROVED/EXACT: no difficulty reduction]\label{c:dichot}
For every $R$, once either $\mathrm{LB}_R$ or $\mathrm{Bridge}_R$ is proved, the other one
is equivalent to $\mathrm{P}\ne\mathrm{NP}$ (or implies it).
Splitting the problem into a lower bound plus a bridge therefore never makes the
second step easier than P vs NP.
\end{corollary}

\begin{theorem}[PROVED/EXACT: the bridge does not relativize for weak $R$]\label{t:rel}
Let $R$ be oracle-free with $\mathrm{LB}_R$ (for example, resolution). Then no relativizing
argument proves $\mathrm{Bridge}_R$.
\end{theorem}
\begin{proof}
$\WR$ does not depend on the oracle. Relative to an oracle $O$ with
$\mathrm{P}^O=\mathrm{NP}^O$ (Baker--Gill--Solovay), the relativized bridge would give
polynomial $\WR$, which contradicts $\mathrm{LB}_R$.
\end{proof}

\begin{theorem}[KNOWN THEOREM: Cook 1975; Kraj\'{\i}\v{c}ek--Pudl\'ak 1989/90]\label{t:ef}
If $A$ runs in polynomial time and $\mathrm{PV}\vdash\forall F\,(A(F)=0\to F\notin\mathrm{SAT})$,
then Extended Frege p-simulates $S_A$. Consequently, ``feasibly provable
$\mathrm{SAT}\in\mathrm{P}$'' implies $\mathrm{Bridge}_{\mathrm{EF}}$.
\end{theorem}
An unconditional $\mathrm{Bridge}_{\mathrm{EF}}$ is OPEN: it would require EF to simulate
algorithms whose correctness is true but not provable in PV. $\mathrm{LB}_{\mathrm{EF}}$ is OPEN
and is the central problem of proof complexity. Whether p-optimal systems exist is OPEN
(Kraj\'{\i}\v{c}ek--Pudl\'ak). They exist if $\mathrm{E}=\mathrm{NE}$, and with advice (M\"uller).

\section{The free optimal sequence clarification}
The directive's point is correct within its scope. Minimizing over derivations is a valid
lower bound for every algorithm whose run \emph{is} an $R$-derivation, because discovering
the sequence only adds work. Examples: DPLL runs are tree-like resolution refutations;
CDCL runs are resolution refutations (Pipatsrisawat--Darwiche 2011).

It is not a lower bound for algorithms whose runs are not $R$-derivations. Theorem~\ref{t:opt}
states this exactly: the transfer needs $R$ to absorb \emph{every} polynomial-time decider.

\paragraph{COMPUTATIONAL hostile control.} The script \texttt{worst-case-derivation/treeres\_php.py}
(0.7\,s total) computes, by exact memoized search with branch-and-bound, the minimum number of
leaves in a tree-like resolution refutation (equivalently, a DPLL tree) of the pigeonhole
formula $\mathrm{PHP}_n^{n+1}$:
\begin{center}
\begin{tabular}{ccccc}
$n$ & vars & clauses & min leaves & min tree size\\\hline
1&2&3&3&5\\ 2&6&9&11&21\\ 3&12&22&43&85
\end{tabular}
\end{center}
The asymptotic trend is the known $2^{\Theta(n\log n)}$ tree-like bound (Iwama--Miyazaki;
Beame--Pitassi), and general resolution is $2^{\Omega(n)}$ (Haken). Meanwhile a decider that
counts $m>n$ uses linear time on this family. So on structured families, a fast
\emph{algorithm} does not induce short derivations in a fixed weak $R$, and $\mathrm{Bridge}_R$
cannot be proved by pointwise simulation of an algorithm into weak $R$.

\section{Barrier-transfer and PQ ledger}
\begin{itemize}
\item Optimization duality (cutting planes / LP): \textbf{failed transfer}. Integrality-gap lower
bounds are $\mathrm{LB}_R$ results for a fixed $R$, so they land under Theorem~\ref{t:collapse}.
\item PQ item 5 (external extraction from a black-box solver): \textbf{analogy/candidate}. It is
exactly the missing ``$R$ absorbs $S_A$'' arrow of Theorem~\ref{t:opt}. Known obstruction:
Kraj\'{\i}\v{c}ek--Pudl\'ak optimality is open.
\item The knot/link handoff of 2026-11-09 is \textbf{failed transfer as a bridge}. A topological
diagram-width lower bound is again an $\mathrm{LB}_R$ for a fixed calculus.
\end{itemize}

\section{Corpus reviewed; duplication / contradiction audit}
Surveyed: the repository tree; \texttt{AGENTS.md}; the root \texttt{main.tex} ownership markers;
the newest five-role handoffs (2026-11-11 LSPV1--5); the prior Claude reports (world ledger,
TPORT, ONETAPE); and a repository grep for $\overline W_R$, ``derivation cost'' and
``worst-case derivation''. No earlier artifact defines $\overline W_R$ or $\mathrm{Bridge}_R$,
so there is no duplication. This run is consistent with world-ledger rows that mark
P=NP-equivalent statements as reformulations: Theorem~\ref{t:collapse} places
$\mathrm{Bridge}_{\mathrm{Res}}$ in exactly that category.
\textbf{What changed relative to root main:} nothing in \texttt{main.tex}, because Claude has
no ownership region there and no edit was made.
\textbf{Adopted:} the free-sequence lower-bound principle, scoped to $R$-trace algorithms.
\textbf{Retired:} ``the weakest sufficient bridge for a resolution-strength $R_{\mathrm{hyp}}$
is weaker than P$\ne$NP'' (FALSIFIED/RETIRED by Theorem~\ref{t:collapse}).
\textbf{Stale or overstated claims:} none found in the files read.

\section{Five-role next targets}
\begin{enumerate}
\item Lead: choose $R$ strong enough that the bridge can be proved. EF under
Theorem~\ref{t:ef}, or a candidate optimal system. State the open $\mathrm{LB}_R$ honestly.
\item Reduction builder: certify that $R_{\mathrm{hyp}}$ is p-equivalent to resolution or to
$\mathrm{Res}(\oplus)$. If it is, Theorem~\ref{t:collapse} applies.
\item Hypergraph certificates: extend the exact search to regular or DAG resolution on small
Tseitin and PHP instances, in chunks of 30\,s or less.
\item Resource auditor: check that $\operatorname{cost}$ (time plus peak space) stays
polynomially related to $|\pi|$. Otherwise Theorem~\ref{t:opt} needs a width or space variant.
\item Red team: look for an $R$ where both $\mathrm{Bridge}_R$ and $\mathrm{LB}_R$ seem provable.
Corollary~\ref{c:dichot} predicts that any such candidate is a proof of P$\ne$NP and needs
extreme scrutiny.
\end{enumerate}

\paragraph{Sources consulted this run.} Wikipedia, ``Propositional proof system'' (equivalence
of polynomial boundedness with NP=coNP; open existence of optimal systems).
M\"uller, ``Optimal proof systems with advice'' (ti1.uni-jena.de).
Kraj\'{\i}\v{c}ek scans, karlin.mff.cuni.cz.
Atserias, STACS 2025 invited talk (drops.dagstuhl.de).

P versus NP remains OPEN.
\end{document}
% END_OF_VERIFIED_REPORT::CLAUDE::WCDERIV-0001
